\section{Discussion}
\label{sec:discussion}

Predicate extraction errors do not affect the final decision equally, and their impact depends on where a predicate sits in the rule structure - invisible to aggregate F1 or accuracy. This supports rule-aware verification: flag only predicates that are both error-prone and high-DIR for the instance's rule path - \textit{is\_written\_assertion} and \textit{is\_introduced\_to\_prove\_truth} in the hearsay setting, and \textit{is\_out\_of\_court} when the assertion clause is otherwise unsatisfied - while deprioritizing low-DIR predicates such as \textit{is\_non\_verbal\_assertive\_conduct}. The taxonomy results add that the same error type carries different risk across models: Gemma 4 31B's hallucinated predicates are decision-critical 96.2\% of the time, versus 46.9-59.5\% for the other two, despite comparable aggregate accuracy - a gap a model-selection decision based on F1 alone would miss.

\textbf{Limitations.} All results come from a single LegalBench task with a shallow rule tree; whether this divergence generalizes to deeper or exception-laden rule structures, and to joint (multi-predicate) intervention, is left to future work. Some DIR estimates rest on as few as 3-5 errors and are directionally informative but statistically fragile, and both the gold table and the rule tree's soundness are validated only within this dataset.
